% Source: https://artofproblemsolving.com/community/c3h2033107_lcm_is_square_of_diference
% Stable identity: aops-topic-2033107-q01
Let \(d=\gcd(a,b)\), \(a=dx\), \(b=dy\) with \(\gcd(x,y)=1\) and \(x>y\ge1\). Then \(\operatorname{lcm}(a,b)=dxy\), so the condition is
\[
dxy=d^2(x-y)^2\iff xy=d(x-y)^2 .
\]
Any prime dividing \(x-y\) and \(x\) would divide \(y\), contradicting \(\gcd(x,y)=1\); similarly for \(y\). So \(\gcd(x-y,xy)=1\). Since \((x-y)^2\) divides \(xy\), we get \(x-y=1\). Then \(d=xy=y(y+1)\). Writing \(k=y\),
\[
a=dx=k(k+1)^2,\qquad b=dy=k^2(k+1).
\]
Conversely, for these values \(\gcd(a,b)=k(k+1)\), so \(\operatorname{lcm}(a,b)=\dfrac{ab}{k(k+1)}=k^2(k+1)^2=(a-b)^2\).

\[
\boxed{(a,b)=\left(k(k+1)^2,\ k^2(k+1)\right),\ k=1,2,3,\dots}
\]
For example \((4,2)\), \((18,12)\), \((48,36)\).
